\section*{الفصل \ref{ch:g4:plants-without-seeds} --- نباتات بلا بذور}

\begin{solution}{exo:g4:plants-without-seeds:1}
\omterm{prop:g4:plants-without-seeds:damp}{الطحالب} \omterm{prop:g4:plants-without-seeds:damp}{والسرخسيات}. وهي تتكاثر \omterm{def:g4:plants-without-seeds:spore}{بالأبواغ} --- وهي حبيبات مجهرية
تُطلق بأعداد هائلة.
\end{solution}

\begin{solution}{exo:g4:plants-without-seeds:2}
\omterm{def:g4:plants-without-seeds:spore}{البوغ} مجهري عارٍ --- لا \omterm{def:g3:animals-through-seasons:active}{مخزن غذاء} فيه ولا \omterm{def:g1:plants-around-us:plant}{نبتة} صغيرة جاهزة؛ أما
\omterm{def:g2:from-seed-to-plant:seed}{البذرة} فتحمل الاثنين داخل غلاف واقٍ. \omterm{def:g4:plants-without-seeds:spore}{والأبواغ} تُصنع بالملايين،
\omterm{def:g2:from-seed-to-plant:seed}{والبذور} بالعشرات.
\end{solution}

\begin{solution}{exo:g4:plants-without-seeds:3}
على الوجه السفلي من السعف: صفوف أنيقة من النقط البنية، كل واحدة
كيس فيه آلاف \omterm{def:g4:plants-without-seeds:spore}{الأبواغ} ينشقّ عند النضج فيطلقها تطير كالغبار.
\end{solution}

\begin{solution}{exo:g4:plants-without-seeds:4}
الطحلب يتكاثر \omterm{def:g4:plants-without-seeds:spore}{بالأبواغ}، \omterm{def:g4:plants-without-seeds:spore}{والأبواغ} ومراحلها الصغيرة الرقيقة تحتاج
إلى الندى --- فهي تجفّ وتموت في الهواء الجاف. والجدار الشمالي
المظلل يبقى نديًّا؛ أما الجدار الجنوبي المشمس فيحترق.
\end{solution}

\begin{solution}{exo:g4:plants-without-seeds:5}
صنع \omterm{def:g1:plants-around-us:plant}{نبتة} جديدة مباشرة من قطعة من \omterm{def:g1:plants-around-us:plant}{النبتة} الأم --- مدّاد، أو \omterm{ex:g4:plants-without-seeds:bulbtuber}{بصلة}،
أو \omterm{ex:g4:plants-without-seeds:bulbtuber}{درنة}، أو عقلة. ووالد واحد فحسب.
\end{solution}

\begin{solution}{exo:g4:plants-without-seeds:6}
ترسل \omterm{def:g1:plants-around-us:plant}{النبتة} ساقًا طويلة رفيعة --- هي المدّاد --- على التربة؛ وحيث
تلامس الأرض تتجذّر وتطلع أوراقًا، فتصير \omterm{def:g1:plants-around-us:plant}{نبتة} صغيرة كاملة؛ ثم يذبل
المدّاد فيقف الابن حرًّا، مطابقًا لوالده.
\end{solution}

\begin{solution}{exo:g4:plants-without-seeds:7}
البطاطا المزروعة (وهي \omterm{ex:g4:plants-without-seeds:bulbtuber}{درنة}) \omterm{def:g2:from-seed-to-plant:germination}{تنبت} من ``عيونها'' على نحو موثوق،
وكل \omterm{def:g1:the-five-senses:sense}{عين} تستطيع أن تنمو \omterm{def:g1:plants-around-us:plant}{نبتة} كاملة مطابقة للصنف الجيد المزروع ---
وهذا أسرع وأوثق من \omterm{def:g2:from-seed-to-plant:seed}{البذور}.
\end{solution}

\begin{solution}{exo:g4:plants-without-seeds:8}
الجواب مفتوح. تظهر \omterm{def:g1:plants-around-us:parts}{جذور} بيضاء من \omterm{def:g1:plants-around-us:parts}{الساق} المغمورة، بعد أسبوع أو
أسبوعين عادة؛ وعندئذ تكون العقلة جاهزة للغرس --- نعناع جديد
كامل، من والد واحد.
\end{solution}

\begin{solution}{exo:g4:plants-without-seeds:9}
\omterm{def:g4:plants-without-seeds:spore}{البوغ} عارٍ رخيص: فمن غير \omterm{def:g3:animals-through-seasons:active}{مخزن غذاء} ولا \omterm{def:g1:plants-around-us:plant}{نبتة} صغيرة يضيع كل واحد
منها تقريبًا، فوجب صنع الملايين. أما \omterm{def:g2:from-seed-to-plant:seed}{البذرة} فعدة نجاة غالية، فتكفي
عشرات منها حسنة التجهيز. وهما صدى لسمكة القدّ (ملايين غير معتنًى
بها) وللفيلة (قليل حسن التزويد).
\end{solution}

\begin{solution}{exo:g4:plants-without-seeds:10}
الفراولة جاءت بالمدّاد --- تكاثرًا خضريًّا، بوالد واحد، ونسخًا
مطابقة. أما الفول فجاء من \omterm{def:g2:from-seed-to-plant:seed}{بذور} صنعتها أزهار، مازجةً إسهامي والدين
(حبوب لقاح من \omterm{def:g3:flowers-fruits-seeds:parts}{زهرة}، \omterm{def:g4:animal-reproduction:sexual}{وبويضات} في أخرى)، فاختلفت كل \omterm{def:g1:plants-around-us:plant}{نبتة} فول اختلافًا
يسيرًا.
\end{solution}

\begin{solution}{exo:g4:plants-without-seeds:11}
\omterm{def:g4:plants-without-seeds:vegetative}{تكاثر خضري} --- \omterm{def:g1:plants-around-us:tree}{فالغصن} عقلة طبيعية. وقد فعل الفيضان عمل مقصّ
البستاني: قطع الفرع، وحمله، وغرسه في طين رطب، وهو المهد المثالي
لتجذّر الصفصاف. والأنهار تصطفّ لنفسها بالصفصاف إذ تكسره وتغرسه في
حركة واحدة.
\end{solution}
